\hdr{§9}{向量代数与空间解析几何：填空 5 分，背公式即可}

\textbf{这一章在竞赛里的样子。}填空题的一道，考平面／直线方程、距离、夹角，或旋转曲面方程。全是公式题，不需要“想”，只需要记准。

\subhdr{9.1 三个乘法}

设 $\boldsymbol a=(a_1,a_2,a_3)$，$\boldsymbol b=(b_1,b_2,b_3)$。

\begin{itemize}[leftmargin=1.6em,itemsep=2pt]
\item \textbf{点积}：$\boldsymbol a\cdot\boldsymbol b=a_1b_1+a_2b_2+a_3b_3=|\boldsymbol a||\boldsymbol b|\cos\theta$。用于求夹角、判断垂直（$\cdot=0$）、求投影。
\item \textbf{叉积}：
\[
\boldsymbol a\times\boldsymbol b=
\begin{vmatrix}
\boldsymbol i&\boldsymbol j&\boldsymbol k\\
a_1&a_2&a_3\\
b_1&b_2&b_3
\end{vmatrix},
\]
$|\boldsymbol a\times\boldsymbol b|$ 是以 $\boldsymbol a,\boldsymbol b$ 为邻边的平行四边形面积；方向由右手法则定。用于求法向量、判断平行（$\times=\boldsymbol0$）。
\item \textbf{混合积}：$(\boldsymbol a,\boldsymbol b,\boldsymbol c)=\boldsymbol a\cdot(\boldsymbol b\times\boldsymbol c)$，绝对值是以三者棱的平行六面体体积。三向量共面 $\iff$ 混合积为 0。
\end{itemize}

\subhdr{9.2 平面与直线}

\begin{center}
\small
\tabtop
\begin{tabular}{@{}p{3.0cm}p{5.4cm}p{4.6cm}@{}}
\thead{对象} & \thead{方程} & \thead{要素}\\
\hline
平面（点法式） & $A(x-x_0)+B(y-y_0)+C(z-z_0)=0$ & 法向量 $\boldsymbol n=(A,B,C)$\\
平面（一般式） & $Ax+By+Cz+D=0$ & 同上\\
直线（对称式） & $\dfrac{x-x_0}{l}=\dfrac{y-y_0}{m}=\dfrac{z-z_0}{n}$ & 方向向量 $\boldsymbol s=(l,m,n)$\\
直线（参数式） & $x=x_0+lt,\ y=y_0+mt,\ z=z_0+nt$ & 同上\\
\end{tabular}
\tabbot
\end{center}

\textbf{距离与夹角}：
\[
d=\frac{|Ax_0+By_0+Cz_0+D|}{\sqrt{A^2+B^2+C^2}}\quad(\text{点到平面}),
\]
\[
\cos\theta=\frac{|\boldsymbol s_1\cdot\boldsymbol s_2|}{|\boldsymbol s_1||\boldsymbol s_2|}\ (\text{线线}),
\qquad
\sin\varphi=\frac{|\boldsymbol s\cdot\boldsymbol n|}{|\boldsymbol s||\boldsymbol n|}\ (\text{线面}).
\]
线面夹角用 $\sin$，因为方向向量与法向量夹角的余角才是线面角。

\textbf{异面直线距离}：$\displaystyle d=\frac{|(\overrightarrow{M_1M_2},\boldsymbol s_1,\boldsymbol s_2)|}{|\boldsymbol s_1\times\boldsymbol s_2|}$，其中 $M_1,M_2$ 是两直线上各一点。

\callout{\textbf{例 1（点到平面）} 求点 $P(1,1,1)$ 到平面 $x+2y+2z-3=0$ 的距离。
\[
d=\frac{|1+2+2-3|}{\sqrt{1+4+4}}=\frac{2}{3}.
\]
\textbf{要点}：分子是“把点代进去取绝对值”，分母是法向量的模。代入前先把平面方程整理成 $Ax+By+Cz+D=0$ 的标准形。}

\subhdr{9.3 曲面：旋转曲面与常见二次曲面}

\textbf{旋转曲面}：$xOy$ 面上的曲线 $F(x,y)=0$ 绕 $x$ 轴旋转，所得曲面方程为
\[
F\bigl(x,\ \pm\sqrt{y^2+z^2}\bigr)=0 .
\]
口诀：\emph{绕谁转，谁不变；另一个变量换成“剩下两个的平方和开根”}。

\callout{\textbf{例 2（旋转曲面）} 曲线 $y^2=2x$（$xOy$ 面上）绕 $x$ 轴旋转，求曲面方程。

绕 $x$ 轴转，$x$ 不变，$y$ 换成 $\pm\sqrt{y^2+z^2}$：
\[
(\pm\sqrt{y^2+z^2})^2=2x\ \Longrightarrow\ y^2+z^2=2x .
\]
这是旋转抛物面。}

\begin{center}
\small
\tabtop
\begin{tabular}{@{}ll@{}}
\thead{名称} & \thead{标准方程}\\
\hline
椭球面 & $\dfrac{x^2}{a^2}+\dfrac{y^2}{b^2}+\dfrac{z^2}{c^2}=1$\\
椭圆抛物面 & $\dfrac{x^2}{a^2}+\dfrac{y^2}{b^2}=z$\\
单叶双曲面 & $\dfrac{x^2}{a^2}+\dfrac{y^2}{b^2}-\dfrac{z^2}{c^2}=1$\\
双叶双曲面 & $\dfrac{x^2}{a^2}-\dfrac{y^2}{b^2}-\dfrac{z^2}{c^2}=1$\\
锥面 & $\dfrac{x^2}{a^2}+\dfrac{y^2}{b^2}=\dfrac{z^2}{c^2}$\\
柱面 & $F(x,y)=0$（母线平行于 $z$ 轴）\\
\end{tabular}
\tabbot
\end{center}

\textbf{识别技巧}：看平方项的系数符号。三个正号且等于 1 是椭球；两个正一个负、右边为 1 是单叶双曲面；一个正两个负、右边为 1 是双叶双曲面；右边是 $z$（一次）是抛物面；右边为 0 是锥面。

\warnbox{\textbf{本章三个高频失分点。}
(1)线面夹角用了 $\cos$（应为 $\sin$）；
(2)旋转曲面忘了 $\pm$，导致只写出一半；
(3)平面方程里 $D$ 的符号代错——先化成标准形再代点。}
